% !TEX root = ../algebra_02.tex

\section{Подгруппа, порождённая подмножеством}

\begin{Definition}{Порождённая подгруппа}{}
	Пусть $M \subset G$. Подгруппа $H_0 < G$ называется подгруппой, порождённой множеством $M$, если выполняются два условия:
	\begin{enumerate}
		\item $M \subset H_0$;
		\item Для любой подгруппы $H < G$, если $M \subset H$, то $H_0 \subset H$ (минимальность).
	\end{enumerate}
	Порождённая подгруппа обозначается $\langle M \rangle$.
\end{Definition}

Эквивалентно, порождённую подгруппу можно определить как пересечение всех подгрупп, содержащих $M$:
\[
\langle M \rangle = \bigcap_{\substack{ H < G \\ M \subset H}} H
\].

Существует явное (конструктивное) описание элементов порожденной подгруппы:
\begin{Proposition}{Явный вид порожденной подгруппы}{}
	Пусть $G$ --- группа и $M \subset G$. Тогда:
	\[
	\langle M \rangle = \{ g_1\cdot g_2 \cdot \ldots \cdot g_n \mid n \ge 0; \ \ g_i \in M \cup M^{-1} \}
	\]
	(При $n=0$ произведение считается равным нейтральному элементу $e$ ).
\end{Proposition}
\begin{proof}
	Обозначим множество всевозможных произведений через $H_0$.
	Оно содержит $M$, так как можно взять произведения длины 1 из элементов $M$.
	Оно является подгруппой, так как содержит $e$, замкнуто относительно умножения (конкатенация произведений дает снова произведение такого же вида), и замкнуто относительно взятия обратного (обратный к произведению равен произведению обратных в обратном порядке, и если $g_i \in M \cup M^{-1}$, то $g_i^{-1} \in M^{-1} \cup M = M \cup M^{-1}$).
	Любая подгруппа $H$, содержащая $M$, обязана содержать $M^{-1}$, а также все конечные произведения элементов из $M \cup M^{-1}$, следовательно $H \supset H_0$.
	Значит, $H_0 = \langle M \rangle$.
\end{proof}
